\bigskip
\bigskip
\subsubsection*{Խնդիրներ և հարցեր թեմա 12-ի վերաբերյալ}

\begin{enumerate}[label=\thesection.\arabic*.]

% 12.1
\item Եթե ունենք $(X, \rho)$ մետրիկ տարածություն, ապա ամեն մի ոչ դատարկ $A \subseteq X$ ենթատարածության վրա կարելի է սահմանել տոպոլոգիա երկու եղանակով:
\par Առաջին դեպքում դիտարկենք $\rho$ մետրիկի $\rho|_A$ սահմանափակումը $A$-ի վրա և վերցնենք $\rho|_A$ մետրիկով որոշված $\tau[\rho|_A]$ մետրիկական տոպոլոգիան:
\par Երկրորդ դեպքում դիտարկենք $X$-ի $\tau[\rho]$ մետրիկական տոպոլոգիան և վերցնենք նրանով մակածված $\tau[\rho]_A$ տոպոլոգիան $A$-ի վրա:
\par Հարց: Նու՞յնն են արդյոք միշտ $\tau[\rho|_A]$ և $\tau[\rho]_A$ տոպոլոգիաները:
\par Լուծում: Պատասխանը գրեթե ակնհայտ դրական է, իսկ ապացույցը կապված է պարզ սահմանումների մակարդակի:

% 12.2
\item Նկարագրեք թվային ուղղի սովորական տոպոլոգիայից ամբողջ թվերի $\mathbb{Z}$ ենթաբազմության վրա մակածված տոպոլոգիան:

% 12.3
\item Ապացուցեք թվային ուղղի ցանկացած երկու
\begin{enumerate}
    \item[ա)] $(a, b)$ և $(c, d)$ տեսքի միջակայքեր սովորական տոպոլոգիայից մակածված տոպոլոգիաներով հոմեոմորֆ տարածություններ են;
    \item[բ)] նույնը $[a, b]$ և $[c, d]$ տեսքի միջակայքերի դեպքում;
    \item[գ)] եթե միջակայքերը տարբեր տեսքերի են՝ մեկը բաց է կամ փակ է, իսկ մյուսը կիսաբաց է, կամ մեկը բաց է, իսկ մյուսը փակ է, ապա ստացված տոպոլոգիական տարածությունները հոմեոմորֆ չեն:
\end{enumerate}

\begin{hint}
Դիտարկեք $f: (a, b) \to (c, d) \quad f(x) = c + \frac{d-c}{b-a}(x-a)$ շարունակական արտապատկերումը և օգտվելով 1, 2 օրինակներից կիրառելի թեորեմ 6-ը թեմա 11-ում:
\end{hint}

% 12.4
\item Դիցուք $\mathbb{R}^2$ կոորդինատային հարթությունում ունենք երկու $AB$ և $CD$ հատվածներ: Դիտարկելով այդ հատվածները որպես $\mathbb{R}^2$ տարածության ենթաբազմություններ $\mathbb{R}^2$-ից նրանց վրա մակածված տոպոլոգիաներով, կառուցեք որևէ $f: AB \to CD$ հոմեոմորֆիզմ:
\par Լուծում: Ներմուծելով $A(x_1, y_1), B(x_2, y_2), C(x_3, y_3), D(x_4, y_4)$ կոորդինատներ, կազմեք $AB$ և $CD$ ուղիղների պարամետրական հավասարումներ նույն $t \in [0, 1]$ տիրույթով և կատարեք պարամետրի անցում:
\par Որոշ դեպքերում պահանջվող հոմեոմորֆիզմը կարելի է կառուցել երկրաչափականորեն, կենտրոնային կամ զուգահեռ պրոյեկցիաներով, ինչպես ցույց է տրված ներքևի օրինակներում:

% 12.5
\item Ստորև պատկերված են ա) շրջանագիծ, բ) օվալ (և իրեն չհատող փակ գիծ), գ) բազմանկյուն՝ որպես միացյալ հարթության ենթաբազմություններ և թեկուզևիս (?) իրեն չհատող տրեֆիլաձև փակ գիծ՝ որպես եռաչափ տարածության ենթաբազմություն:
\par (Images of circle, oval, polygon, trefoil knot)
\par Ապացուցեք, որ դրանք հոմեոմորֆ են իրար:

\begin{hint}
Շրջանագծի վրա վերցրեք $A_1, A_2, ..., A_6$ կետեր, «տրեֆիլաձևի» վրա՝ $B_1, B_2, ..., B_6$ կետեր:
\par Վերցրեք $f_i: A_i A_{i+1} \to B_i B_{i+1}, i=1, 2, ..., 5$ և $f_6: A_6 A_1 \to B_6 B_1$ հոմեոմորֆիզմներ կառուցեք հոմեոմորֆիզմ... (text cuts off) ... տեսնեմ 3-ում:
\end{hint}

% 12.6
\item Դիտարկենք թվային ուղղի $(-1, 1)$ ենթաբազմությունը նրա վրա սովորական տոպոլոգիայից մակածված տոպոլոգիայով: Ապացուցեք, որ այդ տարածությունը հոմեոմորֆ է $(\mathbb{R}, \text{սովոր})$ տարածությանը:

\begin{hint}
Օգտվել թեմա 1-ի օրինակ 9-ում բերված արտապատկերումից:
\end{hint}

\par Ստորև 12.6-12.9 խնդիրներում կորերը և մակերևույթները դիտարկվում են համապատասխանաբար $(\mathbb{R}^2, \text{սովոր})$ և $(\mathbb{R}^3, \text{սովոր})$ տարածություններից մակածված տոպոլոգիաներով: Ապացուցեք.

% 12.7
\item Ցանկացած պարաբոլ հոմեոմորֆ է ուղիղ գծին:

% 12.8
\item Ցանկացած միակտոր հիպերբոլա հոմեոմորֆ է ուղիղ գծին:

% 12.9
\item Ցանկացած երկկտոր հիպերբոլա հոմեոմորֆ է զուգահեռ ուղիղների զույգի:

% 12.10
\item Ցանկացած կիսագնդային մակերևույթ հոմեոմորֆ է հարթությանը:

% 12.11
\item Ապացուցեք. հաշվելիության առաջին և երկրորդ աքսիոմները ժառանգական հատկություններ են: Այսինքն, եթե որևէ տարածություն բավարարում է առաջին (երկրորդ) աքսիոմին, ապա նրա ցանկացած ենթատարածություն նույնպես բավարարում է այդ նույն աքսիոմին:    
\end{enumerate}
